\documentclass[11pt]{article}
\usepackage{amsmath}
\usepackage{amssymb}
\usepackage[dvipsnames]{xcolor}
\usepackage{graphicx}
\usepackage{newpxtext}
\usepackage{eulervm}
\DeclareMathOperator{\diag}{\mathrm{diag}}
\DeclareMathOperator{\cov}{\texttt cov}
\long\def\ans#1{{\color{Blue}{\em #1}}}
\usepackage[margin=1in]{geometry}
\title{Homework 6}
\author{Your name here}
\date{\today{}}

\begin{document}
\maketitle

\section{Expectations, two ways}

\emph{(Adapted from FCML Ex. 2.4)}

\bigskip

\noindent Let $X$ be a random variable with uniform density, $p(x) =
\mathcal{U}(a,b)$ (see FCML section 2.5.1).
\noindent Work out analytically $\mathbf{E}_{p(x)} \left\{ 35 + 3x - 3x^2 +
0.2x^3 + 0.01x^4 \right\}$ for $a=-1$, $b=9$ (show the steps).
Then, compute a sample-based approximation to this expectation by filling out
the necessary code in \texttt{hw.py}.
What is the value you get for the approximation (using all 1000 generated
samples)?
Include the generated figure \texttt{sample\_based\_approximation.pdf} in the
written portion. In the caption for the figure, include a description of the
trend you see in the plot (you learned how to insert figures with captions in
\LaTeX{} in HW0).

\bigskip

\subsection*{Solution}
%%% Answer START %%%
\ans{%
\begin{align*}
\mathbf{E}_{p(x)} \{ 35 + 3x - 3x^2 + 0.2x^3 + 0.01x^4 \} = \int_{-\infty}^{\infty} \left( 35 + 3x - 3x^2 + 0.2x^3 + 0.01x^4 \right) p(x) \, \mathrm{d}x
\end{align*}

With $X \sim U(a,b)$, then its pdf is $p(x) = \frac{1}{b-a}$, and the expectation is:

{\scriptsize
\begin{align*}
\int_{a}^{b} \left( 35 + 3x - 3x^2 + 0.2x^3 + 0.01x^4 \right) \frac{1}{b - a} \, \mathrm{d}x & = \frac{1}{b - a} \int_{a}^{b} \left( 35 + 3x - 3x^2 + 0.2x^3 + 0.01x^4 \right) \, \mathrm{d}x \\
& = \frac{35x + \frac{3}{2}x^2 - \frac{3}{3}x^3 + \frac{0.2}{4}x^4 + \frac{0.01}{5}x^5 \Bigr|_{a}^{b}}{b - a} \\
& = \frac{%
 \left( 35b + \frac{3}{2}b^2 - b^3 + \frac{0.2}{4}b^4 + \frac{0.01}{5}b^5 \right)
 - \left( 35a + \frac{3}{2}a^2 - a^3 + \frac{0.2}{4}a^4 + \frac{0.01}{5}a^5 \right)}
  {b - a}
\end{align*}
}
With $a = -1$ and $b = 9$, the analytically computed expectation is 18.61.
\begin{figure}[htbp]
    \centering
    \includegraphics[width=0.5\textwidth]{sample_based_approximation.pdf}
    \caption{Plot of sampling approximation over 10,000; the red line
    represents the true expected value.}
\end{figure}
}
%%% Answer END %%%

\section{Workload}

How many hours did you spend on this assignment?

\section*{Acknowledgments}

Cite all the people you've worked with on this homework, as well as any other
resources you used apart from the FCML textbook. If you used generative AI
tools (e.g., ChatGPT), please describe how you used them.

If you did not work with anyone else or use generative AI tools, please say so.

\end{document}
